\section{Feature Ablation Study}

\subsection*{Repeat the classification using only feature 1 to 4. What do you obtain?}
When the classification is repeated using only the first four features, the error rates for the models change as follows:
\begin{itemize}
    \item \textbf{Standard MVG:} 8.0\%
    \item \textbf{Naive Bayes:} 7.7\%
    \item \textbf{Tied MVG:} 9.5\%
\end{itemize}

\subsection*{What can we conclude on discarding the last two features?}
We can conclude that discarding the last two features (Features 5 and 6) degrades the overall performance of the classifiers. For the Standard MVG, the error rate increases from 7.0\% to 8.0\%. This indicates that discarding these features results in a net loss of discriminative information.

\subsection*{Despite the inaccuracy of the assumption for these two features, are the Gaussian models still able to extract some useful information?}
Yes. Even though the Gaussian assumption is visually poor for Features 5 and 6 due to their heavy skewness, the models are still capable of extracting useful information from them. The fact that the error rate is lower when these features are included proves that the increase in information outweighs the modeling inaccuracies.

\subsection*{Repeat the classification using only features 1-2 and 3-4. Which model is better in each case?}
\begin{itemize}
    \item \textbf{Case 1 (Features 1-2):} The Standard MVG is significantly better, achieving an error rate of 36.5\%, while the Tied MVG results in 49.4\% (nearly equivalent to a random guess).
    \item \textbf{Case 2 (Features 3-4):} The performance is nearly identical, with the Tied MVG being slightly better at 9.4\% compared to the Standard MVG at 9.5\%.
\end{itemize}

\subsection*{How is this related to the characteristics of the two classifiers?}
The Standard MVG allows for quadratic decision boundaries by using class-specific covariance matrices, whereas the Tied MVG is restricted to linear boundaries. 
In the case of Features 1-2, where classes differ by variance, the quadratic boundary is essential. For Features 3-4, where classes differ by mean and have similar variances, the linear boundary of the Tied model is sufficient and more robust.

\subsection*{Is the tied model effective at all for the first two features? Why?}
No, the tied model is completely ineffective for the first two features. Because the separation in this subspace is driven almost entirely by the difference in class variances, the tied assumption of equal covariance removes the primary source of discriminative information.

\subsection*{And the MVG? And for the second pair of features?}
The MVG is effective for the first pair because it can model class-specific variances. For the second pair (Features 3-4), both the MVG and Tied models are effective; however, the Tied model is slightly more efficient here as the equal-variance assumption holds true, allowing for more stable parameter estimation.